Physics-informed neural networks (PINNs) require a fresh training
run for any change in the residual: a new wavenumber, a new
mobility ratio, a new loss-weight balance. Operator-learning
methods (FNO, DeepONet) amortise across a parametric family but
require solver-generated paired (input, solution) examples
\emph{in their training loop}. We close that gap. LC-PINN is a
single residual-loss training run that produces a continuous
family of PDE solutions over a parameter $\lambda$ --- a physical
scalar coefficient or a loss-weight vector --- without paired
training data, test-time adaptation, or per-$\lambda$ retuning.
Reference solutions are used only to compute test-time error. The construction adapts loss-conditional
networks \citep{dosovitskiy2020} to PINN training: $\lambda$
becomes a network input drawn fresh from a fixed prior at every
optimisation step. We prove a $\lambda$-invariance result --- at
a global optimum of the $\lambda$-expected loss the network is a
per-$\lambda$ residual minimiser for $p_\lambda$-almost every
$\lambda$ --- with a finite-capacity corollary that bounds the
suboptimality at every $\lambda$ in the support. Across 1D, 2D, and 3D parametric Helmholtz, 1D stationary Schrödinger
with a multiplicative parametric potential, and loss-weight
amortisation on viscous Burgers and Buckley--Leverett, LC reaches
rel-$L^2$ within a factor of two of the strongest single-$\lambda$
baseline. On 1D Helmholtz it is the only method whose rel-$L^2$
stays within one order of magnitude across $k\!\in\![1,10]$ ---
per-$k$ retrained baselines vary by three.
